\section{The Leray spectral sequence}
\label{section-Leray}

\begin{lemma}
\label{lemma-before-Leray}
Let $f : X \to Y$ be a morphism of ringed spaces.
There is a commutative diagram
$$
\xymatrix{
D^{+}(X) \ar[rr]_-{R\Gamma(X, -)} \ar[d]_{Rf_*} & &
D^{+}(\mathcal{O}_X(X)) \ar[d]^{\text{restriction}} \\
D^{+}(Y) \ar[rr]^-{R\Gamma(Y, -)} & &
D^{+}(\mathcal{O}_Y(Y))
}
$$
More generally for any $V \subset Y$ open and $U = f^{-1}(V)$ there
is a commutative diagram
$$
\xymatrix{
D^{+}(X) \ar[rr]_-{R\Gamma(U, -)} \ar[d]_{Rf_*} & &
D^{+}(\mathcal{O}_X(U)) \ar[d]^{\text{restriction}} \\
D^{+}(Y) \ar[rr]^-{R\Gamma(V, -)} & &
D^{+}(\mathcal{O}_Y(V))
}
$$
See also Remark \ref{remark-elucidate-lemma} for more explanation.
\end{lemma}

\begin{proof}
Let
$\Gamma_{res} : \textit{Mod}(\mathcal{O}_X) \to \text{Mod}_{\mathcal{O}_Y(Y)}$
be the functor which associates to an $\mathcal{O}_X$-module $\mathcal{F}$
the global sections of $\mathcal{F}$ viewed as an $\mathcal{O}_Y(Y)$-module
via the map $f^\sharp : \mathcal{O}_Y(Y) \to \mathcal{O}_X(X)$. Let
$restriction : \text{Mod}_{\mathcal{O}_X(X)} \to \text{Mod}_{\mathcal{O}_Y(Y)}$
be the restriction functor induced by
$f^\sharp : \mathcal{O}_Y(Y) \to \mathcal{O}_X(X)$. Note that $restriction$
is exact so that
its right derived functor is computed by simply applying the restriction
functor, see
Derived Categories, Lemma \ref{derived-lemma-right-derived-exact-functor}.
It is clear that
$$
\Gamma_{res}
=
restriction \circ \Gamma(X, -)
=
\Gamma(Y, -) \circ f_*
$$
We claim that
Derived Categories, Lemma \ref{derived-lemma-compose-derived-functors}
applies to both compositions. For the first this is clear by our remarks
above. For the second, it follows from
Lemma \ref{lemma-pushforward-injective} which implies that
injective $\mathcal{O}_X$-modules are mapped to $\Gamma(Y, -)$-acyclic
sheaves on $Y$.
\end{proof}

\begin{remark}
\label{remark-elucidate-lemma}
Here is a down-to-earth explanation of the meaning of
Lemma \ref{lemma-before-Leray}. It says that given
$f : X \to Y$ and $\mathcal{F} \in \textit{Mod}(\mathcal{O}_X)$
and given an injective resolution $\mathcal{F} \to \mathcal{I}^\bullet$
we have
$$
\begin{matrix}
R\Gamma(X, \mathcal{F}) & \text{is represented by} &
\Gamma(X, \mathcal{I}^\bullet) \\
Rf_*\mathcal{F} & \text{is represented by} & f_*\mathcal{I}^\bullet \\
R\Gamma(Y, Rf_*\mathcal{F}) & \text{is represented by} &
\Gamma(Y, f_*\mathcal{I}^\bullet)
\end{matrix}
$$
the last fact coming from Leray's acyclicity lemma
(Derived Categories, Lemma \ref{derived-lemma-leray-acyclicity})
and Lemma \ref{lemma-pushforward-injective}.
Finally, it combines this with the trivial observation that
$$
\Gamma(X, \mathcal{I}^\bullet)
=
\Gamma(Y, f_*\mathcal{I}^\bullet).
$$
to arrive at the commutativity of the diagram of the lemma.
\end{remark}

\begin{lemma}
\label{lemma-modules-abelian}
Let $X$ be a ringed space.
Let $\mathcal{F}$ be an $\mathcal{O}_X$-module.
\begin{enumerate}
\item The cohomology groups $H^i(U, \mathcal{F})$ for $U \subset X$ open
of $\mathcal{F}$ computed as an $\mathcal{O}_X$-module, or computed as an
abelian sheaf are identical.
\item Let $f : X \to Y$ be a morphism of ringed spaces.
The higher direct images $R^if_*\mathcal{F}$ of $\mathcal{F}$
computed as an $\mathcal{O}_X$-module, or computed as an abelian sheaf
are identical.
\end{enumerate}
There are similar statements in the case of bounded below
complexes of $\mathcal{O}_X$-modules.
\end{lemma}

\begin{proof}
Consider the morphism of ringed spaces
$(X, \mathcal{O}_X) \to (X, \underline{\mathbf{Z}}_X)$ given
by the identity on the underlying topological space and by
the unique map of sheaves of rings
$\underline{\mathbf{Z}}_X \to \mathcal{O}_X$.
Let $\mathcal{F}$ be an $\mathcal{O}_X$-module.
Denote $\mathcal{F}_{ab}$ the same sheaf seen as an
$\underline{\mathbf{Z}}_X$-module, i.e., seen as a sheaf of
abelian groups. Let
$\mathcal{F} \to \mathcal{I}^\bullet$ be an injective resolution.
By Remark \ref{remark-elucidate-lemma} we see that
$\Gamma(X, \mathcal{I}^\bullet)$ computes both
$R\Gamma(X, \mathcal{F})$ and $R\Gamma(X, \mathcal{F}_{ab})$.
This proves (1).

\medskip\noindent
To prove (2) we use (1) and Lemma \ref{lemma-describe-higher-direct-images}.
The result follows immediately.
\end{proof}

\begin{lemma}[Leray spectral sequence]
\label{lemma-Leray}
Let $f : X \to Y$ be a morphism of ringed spaces.
Let $\mathcal{F}^\bullet$ be
a bounded below complex of $\mathcal{O}_X$-modules.
There is a spectral sequence
$$
E_2^{p, q} = H^p(Y, R^qf_*(\mathcal{F}^\bullet))
$$
converging to $H^{p + q}(X, \mathcal{F}^\bullet)$.
\end{lemma}

\begin{proof}
This is just the Grothendieck spectral sequence
Derived Categories, Lemma \ref{derived-lemma-grothendieck-spectral-sequence}
coming from the composition of functors
$\Gamma_{res} = \Gamma(Y, -) \circ f_*$ where $\Gamma_{res}$ is as
in the proof of Lemma \ref{lemma-before-Leray}.
To see that the assumptions of
Derived Categories, Lemma \ref{derived-lemma-grothendieck-spectral-sequence}
are satisfied, see the proof of Lemma \ref{lemma-before-Leray} or
Remark \ref{remark-elucidate-lemma}.
\end{proof}

\begin{remark}
\label{remark-Leray-ss-more-structure}
The Leray spectral sequence, the way we proved it in Lemma \ref{lemma-Leray}
is a spectral sequence of $\Gamma(Y, \mathcal{O}_Y)$-modules. However, it
is quite easy to see that it is in fact a spectral sequence of
$\Gamma(X, \mathcal{O}_X)$-modules. For example $f$ gives rise to
a morphism of ringed spaces
$f' :  (X, \mathcal{O}_X) \to (Y, f_*\mathcal{O}_X)$.
By Lemma \ref{lemma-modules-abelian} the terms $E_r^{p, q}$ of the
Leray spectral sequence for an $\mathcal{O}_X$-module $\mathcal{F}$
and $f$ are identical with those for $\mathcal{F}$ and $f'$
at least for $r \geq 2$. Namely, they both agree with the terms of the Leray
spectral sequence for $\mathcal{F}$ as an abelian sheaf.
And since $(f_*\mathcal{O}_X)(Y) = \mathcal{O}_X(X)$ we see the result.
It is often the case
that the Leray spectral sequence carries additional structure.
\end{remark}

\begin{lemma}
\label{lemma-apply-Leray}
Let $f : X \to Y$ be a morphism of ringed spaces.
Let $\mathcal{F}$ be an $\mathcal{O}_X$-module.
\begin{enumerate}
\item If $R^qf_*\mathcal{F} = 0$ for $q > 0$, then
$H^p(X, \mathcal{F}) = H^p(Y, f_*\mathcal{F})$ for all $p$.
\item If $H^p(Y, R^qf_*\mathcal{F}) = 0$ for all $q$ and $p > 0$, then
$H^q(X, \mathcal{F}) = H^0(Y, R^qf_*\mathcal{F})$ for all $q$.
\end{enumerate}
\end{lemma}

\begin{proof}
These are two simple conditions that force the Leray spectral sequence to
degenerate at $E_2$. You can also prove these facts directly (without using
the spectral sequence) which is a good exercise in cohomology of sheaves.
\end{proof}

\begin{lemma}
\label{lemma-higher-direct-images-compose}
\begin{slogan}
The total derived functor of a composition is the
composition of the total derived functors.
\end{slogan}
Let $f : X \to Y$ and $g : Y \to Z$ be morphisms of ringed spaces.
In this case $Rg_* \circ Rf_* = R(g \circ f)_*$ as functors
from $D^{+}(X) \to D^{+}(Z)$.
\end{lemma}

\begin{proof}
We are going to apply
Derived Categories, Lemma \ref{derived-lemma-compose-derived-functors}.
It is clear that $g_* \circ f_* = (g \circ f)_*$, see
Sheaves, Lemma \ref{sheaves-lemma-pushforward-composition}.
It remains to show that $f_*\mathcal{I}$ is $g_*$-acyclic.
This follows from Lemma \ref{lemma-pushforward-injective}
and the description of the
higher direct images $R^ig_*$ in
Lemma \ref{lemma-describe-higher-direct-images}.
\end{proof}

\begin{lemma}[Relative Leray spectral sequence]
\label{lemma-relative-Leray}
Let $f : X \to Y$ and $g : Y \to Z$ be morphisms of ringed spaces.
Let $\mathcal{F}$ be an $\mathcal{O}_X$-module.
There is a spectral sequence with
$$
E_2^{p, q} = R^pg_*(R^qf_*\mathcal{F})
$$
converging to $R^{p + q}(g \circ f)_*\mathcal{F}$.
This spectral sequence is functorial in $\mathcal{F}$, and there
is a version for bounded below complexes of $\mathcal{O}_X$-modules.
\end{lemma}

\begin{proof}
This is a Grothendieck spectral sequence for composition of functors
and follows from Lemma \ref{lemma-higher-direct-images-compose} and
Derived Categories, Lemma \ref{derived-lemma-grothendieck-spectral-sequence}.
\end{proof}














